% ----------------------
\subsection{Section 54 (10 June 1831)}
% ----------------------
	\label{DC54}
	
	% -----
	\subsubsection{Historical Background}
	% -----
		\label{DC54-HB}
		
		% --
		\paragraph{JSP}
		% --
		
			``This revelation provided instructions to the church members from Colesville, New York, after they encountered difficulties settling in Thompson, Ohio. In May 1831, shortly after they arrived in Ohio in compliance with revelations directing all New York members to gather there, JS instructed Bishop Edward Partridge to settle the Colesville members in Thompson on land offered by convert Leman Copley, a former Shaker.
			
			``When JS moved to Ohio in early February, Copley had invited JS and Sidney Rigdon to live with him at Thompson, offering to `furnish them houses \& provisions \&c.' While both JS and Rigdon made other living arrangements for their families, Joseph Knight Sr., who accompanied JS to Ohio, recalled that in March he and JS went to Thompson, presumably to see about settling the soon-to-be-emigrating Colesville congregation on Copley's property. Once the Colesville members arrived in Ohio, JS sent them to Thompson to live on Copley's extensive landholdings.
			
			``The arrangement with Copley apparently granted the Colesville members the privilege to live on the property in return for making improvements upon it, and according to Joseph Knight Sr., they `all went to work and made fence and planted and sowed the fields.' On  May Copley was called to preach, along with Sidney Rigdon and Parley P. Pratt, to the Shaker settlement in nearby North Union, Ohio.7 However, the missionary expedition failed to convert any of the Shakers, and the resulting confrontation between the Mormon elders and the Shakers apparently disturbed Copley. He soon went back to North Union, where he apparently reconciled with the Shaker community, and he then returned to Thompson with Shaker leader Ashbel Kitchell, perhaps intending to evict the Mormons. During his visit to Thompson, Kitchell held a meeting with the Mormons on Copley's farm, was involved in a contentious altercation, and initiated efforts to remove them. Joseph Knight Jr. recalled, `We had to leave his [Copley's] farm and pay sixty dollars damage,' adding bitterly that the payment was for `fitting up his houses and planting his ground.'
			
			``Because of the difficulties with Copley, Newel Knight, the presiding elder over the Colesville group, went to Kirtland to consult with JS before the conference held in early June. Knight later explained that as a result of a revelation on 6 June, the last day of that conference, `we now understood that this [Ohio] was not the land of our inheritance---the land of promise, for it was made known in a revelation, that Missouri was the place chosen for the gathering of the Church, and several were called to lead the way to that state.' Though that 6 June revelation addressed the church generally, the 10 June revelation responded specifically to the concerns of the Colesville members living in Thompson. Knight later introduced the revelation with these words: `As I had come to see brother Joseph concerning our position in Thompson, he enquired of the Lord and received the following revelation.' John Whitmer similarly recalled, `At this time the Church at Thompson Ohio was involved in difficulty, becaus of the rebellion of Leman Copley. Who would not do as he had previously agreed. Which thing confused the whole church and finally the Lord spake unto Joseph Smith Jr the prophit.'''
			
		% --
		\paragraph{HDDC}
		% --
		
			HDDC 701--702 has a long entry from Jared Carer's journal describing the ``fighting and bickering among the people,'' which also made this situation more difficult.
		
	% -----
	\subsubsection{Versions}
	% -----
		\label{DC54-V}
		
		See the \href{https://drive.google.com/file/d/1goAvNz8jS4R8slYfMG_db55Ty2z_vmgK/view?usp=sharing}{sections comparison} document.
	
		\begin{compactdesc}
			\item[JSP] \hyperref[EMS-1831-JS-10-JUN-1831]{10 June 1831}
			\item[BC] Chapter 56, 128--129
			\item[DC 1835] Section 67, 195--196
			\item[TS] 5:4 (15 February 1844), 432
			\item[DC 1844] Section 68, 297--298
			\item[MS] 5:8 (January 1845), 113
		\end{compactdesc}

	% -----
	\subsubsection{Section 54:1} % *DC54-1 
	% -----
		\label{DC54-1}

		BEHOLD, thus saith the Lord, even Alpha and Omega,%
			\footnote{%
				% \label{}%
				See \hyperref[REVELATIONNT-1-8]{Revelation 1:8} and notes there.
				}
		the beginning and the end, even he who was crucified for the sins of the world---

		% \begin{framed}
		% \scriptsize
			% \begin{compactdesc}
				% \item[JSP] 
				% % \item[BC] 
				% % \item[]
			% \end{compactdesc}
		% \end{framed}

	% -----
	\subsubsection{Section 54:2} % *DC54-2 
	% -----
		\label{DC54-2}

		Behold, verily, verily, I say unto you, my servant Newel Knight, you shall stand fast in the office whereunto I have appointed you.

		% \begin{framed}
		% \scriptsize
			% \begin{compactdesc}
				% \item[JSP] 
				% % \item[BC] 
				% % \item[]
			% \end{compactdesc}
		% \end{framed}

	% -----
	\subsubsection{Section 54:3} % *DC54-3 
	% -----
		\label{DC54-3}

		And if your brethren desire to escape their enemies,%
			\footnote{%
				% \label{}%
				Compare the \hyperref[SALVATION-JS]{much later JS remark}, from May 1843, about salvation being to be saved or possibly ``severed'' from all your enemies.
				}
		let them repent of all their sins, and become truly \hyperref[HUMILITY]{humble} before me and \hyperref[CONTRITE]{contrite}.

		% \begin{framed}
		% \scriptsize
			% \begin{compactdesc}
				% \item[JSP] 
				% % \item[BC] 
				% % \item[]
			% \end{compactdesc}
		% \end{framed}

	% -----
	\subsubsection{Section 54:4} % *DC54-4 
	% -----
		\label{DC54-4}

		And as the covenant which they made unto me has been broken, even so it has become void and of none effect.%
			\footnote{%
				% \label{}%
				The way the Lord puts it here makes it seem as though any ``punishment'' associated with the covenant is also ``void and of none effect.''
				}

		% \begin{framed}
		% \scriptsize
			% \begin{compactdesc}
				% \item[JSP] 
				% % \item[BC] 
				% % \item[]
			% \end{compactdesc}
		% \end{framed}

	% -----
	\subsubsection{Section 54:5} % *DC54-5 
	% -----
		\label{DC54-5}

		And wo to him by whom this \hyperref[OFFENSE]{offense} cometh, for it had been better for him that he had been drowned in the depth of the sea.

		% \begin{framed}
		% \scriptsize
			% \begin{compactdesc}
				% \item[JSP] 
				% % \item[BC] 
				% % \item[]
			% \end{compactdesc}
		% \end{framed}

	% -----
	\subsubsection{Section 54:6} % *DC54-6 
	% -----
		\label{DC54-6}

		But blessed are they who have kept the covenant and observed the commandment, for they shall obtain mercy.%
			\footnote{%
				% \label{}%
				This language reminds me a bit of the Beatitudes and \hyperref[3NEPHI-12-7]{3 Nephi 12:7}.
				}

		% \begin{framed}
		% \scriptsize
			% \begin{compactdesc}
				% \item[JSP] 
				% % \item[BC] 
				% % \item[]
			% \end{compactdesc}
		% \end{framed}

	% -----
	\subsubsection{Section 54:7} % *DC54-7 
	% -----
		\label{DC54-7}

		Wherefore, go to now and flee the land, lest your enemies come upon you; and take your journey, and appoint whom you will to be your leader, and to pay moneys for you.

		% \begin{framed}
		% \scriptsize
			% \begin{compactdesc}
				% \item[JSP] 
				% % \item[BC] 
				% % \item[]
			% \end{compactdesc}
		% \end{framed}

	% -----
	\subsubsection{Section 54:8} % *DC54-8 
	% -----
		\label{DC54-8}

		And thus you shall take your journey into the regions westward, unto the land of Missouri, unto the borders of the Lamanites.

		% \begin{framed}
		% \scriptsize
			% \begin{compactdesc}
				% \item[JSP] 
				% % \item[BC] 
				% % \item[]
			% \end{compactdesc}
		% \end{framed}

	% -----
	\subsubsection{Section 54:9} % *DC54-9 
	% -----
		\label{DC54-9}

		And after you have done journeying, behold, I say unto you, seek ye a living like unto men, until I prepare a place for you.

		% \begin{framed}
		% \scriptsize
			% \begin{compactdesc}
				% \item[JSP] 
				% % \item[BC] 
				% % \item[]
			% \end{compactdesc}
		% \end{framed}

	% -----
	\subsubsection{Section 54:10} % *DC54-10 
	% -----
		\label{DC54-10}

		And again, be patient in tribulation until I come; and, behold, I come quickly, and my reward is with me, and they who have sought me early shall find rest to their souls.%
			\footnote{%
				% \label{}%
				Lots of scriptures to point to here, but you could check out \hyperref[MORONI-9-6]{Moroni 9:6} and \hyperref[DC84-24]{DC 84:24}.
				}
		Even so.  Amen.

		% \begin{framed}
		% \scriptsize
			% \begin{compactdesc}
				% \item[JSP] 
				% % \item[BC] 
				% % \item[]
			% \end{compactdesc}
		% \end{framed}